\section{为什么把治理与会计放入投资结论}

对周期性制造企业，利润预告的质量不仅取决于损益表，还取决于收入确认、存货成本、资产减值、经营性现金流和关联交易等口径能否相互印证。雅化在 2025A 与 2026Q1 已出现利润修复而经营现金为负、应收和存货增加的组合；这不自动说明会计质量存在问题，却使正式 H1 的附注、现金流量表和营运资本解释成为模型的必要输入。评级和目标价已对这一不确定性施加概率折价，而非在附录中才披露。

\begin{exhibitbox}[表 A-9：正式 H1 的会计与治理检查清单]
\begin{tabularx}{\exhibitboxwidth}{L{2.9cm}X X}
\toprule
检查主题 & 要核对的披露与勾稽 & 对估值的后果 \\
\midrule
收入确认与合同资产 & 锂盐、民爆产品和爆破服务收入是否与应收、合同资产及回款方向一致 & 若收入增长不能转为回款，维持现金折价 \\
存货与跌价准备 & 原料、在产品、产成品的金额、周转、跌价准备及价格变化解释 & 若存货上升且价格回落，重新评估毛利与减值风险 \\
经营现金流 & 净利润与 OCF 的差额、经营性应收/存货变动、税费支付 & OCF 是基准 PE 能否上修的先决条件 \\
资产减值与投资收益 & 减值、金融资产公允价值、处置与其他非经营项目的来源 & 将一次性项目从 H2 常态利润中剔除 \\
有息负债与资金用途 & 短借、长借、受限资金、资本开支与项目付款计划 & 净现金代理不再成立时，调整 EV 交叉检查 \\
关联方与少数股东 & 关联采购/销售、少数股东损益和资金往来 & 防止将并表利润全部归属母公司 \\
民爆合规与安全 & 许可证、产量、项目执行、安全生产与环保披露 & 许可或项目变化不能脱离实际产量估值 \\
\bottomrule
\end{tabularx}
\end{exhibitbox}

\section{披露控制：每条可进入模型的结论都需要可追溯来源}

本报告的来源分为四层。第一层是公司公告、年报、季报和权益分派资料，用于历史财务、预告区间、股本和许可等可核验事实；第二层是完整券商 PDF，用于外部盈利假设和目标价差异；第三层是行业价格、下游销量和同行估值，用于周期与相对估值交叉检查；第四层是关系、项目、客户或设计能力线索，只能进入监控表。来源登记表记录了文件路径、日期、用途和质量等级，主张审计表则把每个材料结论连到其估值后果。

\begin{exhibitbox}[表 A-10：证据到模型的控制链]
\begin{tabularx}{\exhibitboxwidth}{L{2.6cm}L{3.0cm}X X}
\toprule
证据层级 & 代表材料 & 可以支持的结论 & 不能支持的结论 \\
\midrule
公司一级披露 & 2025 年报、2026Q1、H1 预告 & 历史收入/毛利/现金、H1 利润区间、股本 & 未披露的 H2 价格、订单、成本和现金 \\
完整卖方报告 & 国信、开源、东吴 PDF & 外部 EPS、方法和目标价分歧 & House EPS 或目标价的直接输入 \\
行业与可比资料 & 商品价格、NEV 销量、同行 PE & 周期位置、需求方向、估值区间校验 & 雅化的实际实现价、份额或订单 \\
关系与线索资料 & 资源权益、许可、项目和客户链 & 战略观察、上行情景验证条件 & 基准收入、利润和独立资产价值 \\
\bottomrule
\end{tabularx}
\end{exhibitbox}

\section{分析师独立性与版本控制}

截至 2026-07-23，市场价格、估值模型、摘要目标和 PDF 统一使用 17.75 元；任何仍含 16.79 元的旧事件稿文件均未被主文引用。外部卖方的盈利预测和目标价与 House 模型分开存储，最终目标中的外部权重为零。模型复算通过并不改变其“中低置信度”的事实：它只表示输入、公式和四舍五入是一致的。正式 H1 发布后应创建新的资料截点，保留本版本作为事后检验基准，而不是覆盖历史版本。

\begin{sourcequalitybox}[本报告的责任边界]
研究结论为“减持 / 避免新增仓位”，原因是现价接近概率加权价值且关键经营验证尚未完成；并非对公司长期战略或行业价格作绝对否定。若后续披露改变现金、单位经济或 H2 归母的可验证范围，评级、目标和情景概率必须同步重算。
\end{sourcequalitybox}

\sourcenote{来源登记册、主张审计、检索边界日志与估值复算工作底稿。}
